\section{Related Work}
\label{sec:related}

Quantitative evaluation of board games has been approached from several directions.
BoardGameGeek's community-assigned weight scores provide the most widely-used
complexity proxy~\cite{woods2012eurogames}, but weight is a single dimension and
conflates mechanical complexity with decision complexity. Zagal et al.'s Ludic
Ontology~\cite{zagal2008collaborative} provides a richer vocabulary for game
components but was not designed for quantitative scoring. Elverdam and Aarseth's
game classification framework~\cite{elverdam2007game} identifies structural
dimensions but lacks empirical validation on a large corpus.

In adjacent domains, evaluation rubrics for digital games have been studied by
Susi and Johannesson~\cite{susi2000play} and Sweetser and Wyeth~\cite{sweetser2005gameflow},
but the mechanics of physical tabletop games differ substantially from video games.
The TIGRIS corpus design draws most directly on psychometric traditions in educational
measurement~\cite{cronbach1951coefficient}: treating axis scores as items in a
multi-dimensional instrument and establishing reliability via Cronbach's $\alpha$.
